% ============================================================
% github-ssh.tex — A SUBDOCUMENT
% This file can be compiled two different ways:
%   1. On its own: run "pdflatex github-ssh.tex" directly.
%      It will borrow main.tex's preamble (packages/styles)
%      and produce a small standalone PDF of just this section.
%   2. As part of the whole guide: when main.tex runs
%      \subfile{github-ssh}, only the content between
%      \begin{document} and \end{document} below gets inserted
%      into main.tex — the \documentclass line here is ignored.
% ============================================================

\documentclass[main.tex]{subfiles}
% This is NOT a normal \documentclass line. The [main.tex] part
% tells the "subfiles" package: "borrow your preamble settings
% from main.tex." Both files must be in the same folder for
% this to work (or adjust the path, e.g. [../main.tex]).

\begin{document}

\section{Adding an SSH Key to GitHub and Cloning a Repository}
% Top-level heading for this content. If this file is inserted
% into main.tex, this becomes a numbered section there too.

\subsection{Check for Existing SSH Keys}
Open a git bash terminal and check whether you already have an SSH key pair.

\begin{lstlisting}
ls -al ~/.ssh
\end{lstlisting}
% \begin{lstlisting} ... \end{lstlisting} displays text exactly
% as typed, in a monospaced font — perfect for terminal commands.
% It uses the "terminal" style we defined in main.tex's preamble.

Look for files named \texttt{id\_ed25519} / \texttt{id\_ed25519.pub}
or \texttt{id\_rsa} / \texttt{id\_rsa.pub}. If these exist, skip to
the ``Add the Key to GitHub'' step.
% \texttt{...} renders text in a monospace/typewriter font,
% useful for filenames within a normal sentence.
% Note: underscores must be escaped as \_ in LaTeX, since a bare
% underscore normally means "subscript."

\subsection{Generate a New SSH Key}
\begin{lstlisting}
ssh-keygen -t ed25519 -C "your_email@example.com"
\end{lstlisting}

\subsection{Start the SSH Agent and Add Your Key}
\begin{lstlisting}
eval "$(ssh-agent -s)"
ssh-add ~/.ssh/id_ed25519
\end{lstlisting}

\subsection{Copy the Public Key}
\begin{lstlisting}
cat ~/.ssh/id_ed25519.pub
\end{lstlisting}

\subsection{Add the Key to GitHub}
\begin{enumerate}[leftmargin=*]
% enumerate = numbered list. [leftmargin=*] (from the enumitem
% package) removes extra indentation so it aligns neatly.
    \item Log in to GitHub and go to \textbf{Settings}.
    \item Click \textbf{SSH and GPG keys} in the sidebar.
    \item Click \textbf{New SSH key}, give it a title, and paste the key.
    \item Click \textbf{Add SSH key}.
\end{enumerate}
% \textbf{...} makes text bold.

\subsection{Test the Connection}
\begin{lstlisting}
ssh -T git@github.com
\end{lstlisting}

\subsection{Clone a Repository Using SSH}
\begin{lstlisting}
git clone git@github.com:username/repository.git
\end{lstlisting}

\end{document}
% This \end{document} closes the subfile's own document body.
% When compiled standalone, this marks the end of the PDF.
% When pulled in via \subfile{github-ssh} from main.tex, this
% line (and \begin{document}/\documentclass above) is simply
% skipped by the subfiles package — only the content in between
% gets inserted.